\ProvidesPackage{pre}


% -----------------------------------------------------------
% PAGE GEOMETRY AND FONTS
% -----------------------------------------------------------


\usepackage{times} % Serif font
\usepackage[lite,subscriptcorrection,slantedGreek,nofontinfo]{mtpro2}
\pdfmapfile{=mtpro2.map}

% -----------------------------------------------------------
% ESSENTIAL PACKAGES
% -----------------------------------------------------------
\usepackage[english]{babel}
\usepackage{amsmath, amssymb, amsthm, mathtools, yhmath}
\usepackage{graphicx}
\usepackage{tikz}
\usetikzlibrary{calc, patterns}
\usepackage{pgfplots}
\pgfplotsset{compat=newest}
\usepgfplotslibrary{fillbetween}
\usepackage{hyperref}
\usepackage{color, xcolor}
\usepackage{titlesec}
\usepackage{setspace}
\usepackage{fancyhdr}
\pagestyle{fancy}
\usepackage{enumitem}
\usepackage{float}
\usepackage{caption}
\usepackage{subcaption}
\usepackage{multicol}
\usepackage{afterpage}
\usepackage{verbatim}
\usepackage{makecell}
\usepackage{varwidth}
\usepackage{ifthen}
\usepackage{cancel}
\usepackage{needspace}
\usepackage{etoolbox}
\usepackage{xparse}
\usepackage{xstring}
\usepackage{framed}
\usepackage{todonotes}
\usepackage[table]{xcolor}
\usepackage{fontawesome}
\usepackage{soul}


\fancyhf{} % Clear all header/footer fields

% Book-style headers: odd pages right, even pages left
\fancyhead[RO]{\rightmark} % Section name on odd pages
\fancyhead[LE]{\leftmark}  % Chapter name on even pages
\fancyfoot[C]{\thepage}    % Page number in center





% -----------------------------------------------------------
% MATHEMATICS & SYMBOLS
% -----------------------------------------------------------
\usepackage{mathrsfs}
\usepackage{polynom}
\usepackage{tkz-euclide}

% -----------------------------------------------------------
% STRUCTURE AND THEOREMS
% -----------------------------------------------------------
\newtheorem{theorem}{Theorem}[section]
\newtheorem{definition}[theorem]{Definition}
\newtheorem{example}[theorem]{Example}
\newtheorem{prop}[theorem]{Proposition}
\newtheorem{lemma}[theorem]{Lemma}
\newtheorem{axiom}[theorem]{Axiom}
\newtheorem{coro}[theorem]{Corollary}
\newtheorem{remark}[theorem]{Remark}
\renewcommand{\theequation}{\thesection.\arabic{equation}}

% -----------------------------------------------------------
% CUSTOM FORMATTING
% -----------------------------------------------------------
\titleformat{\chapter}[display]
{\normalfont\Large\bfseries\centering}
{\vspace{-2em}\textsc{Chapter \thechapter}}
{1ex}
{\Huge}
[\vspace{0.5ex}\rule{\textwidth}{0.4pt}]

\definecolor{winered}{RGB}{190, 47, 55}
\titleformat{\section}
{\color{winered}\LARGE}
{\thesection}{1em}{}
\renewcommand{\thesection}{\arabic{section}}

\titleformat{\subsection}
{\large\bfseries\itshape}
{\thesubsection}{1em}{}
\renewcommand{\thesubsection}{\thesection.\arabic{subsection}}

\pretocmd{\section}{\needspace{3\baselineskip}}{}{}
\pretocmd{\subsection}{\needspace{2\baselineskip}}{}{}

% -----------------------------------------------------------
% CUSTOM BOXES AND COLORS
% -----------------------------------------------------------
\usepackage[most]{tcolorbox}
\tcbuselibrary{skins, breakable}


\newtcolorbox{fancyhighlightorange}[1][]{
	enhanced,
	breakable,
	colback=white,
	boxrule=0pt,
	frame hidden,
	interior hidden,
	sharp corners,
	left=5pt, right=5pt, top=5pt, bottom=5pt,
	before skip=10pt, after skip=10pt,
	overlay={
		\path [left color=white, right color=white, middle color=orange!10]
		([xshift=-5pt]frame.south west) rectangle ([xshift=5pt]frame.north east);
		\draw [line width=0.4pt, left color=white, right color=white, middle color=orange!60!brown!60]
		([xshift=-5pt]frame.north west) -- ([xshift=5pt]frame.north east);
		\draw [line width=0.4pt, left color=white, right color=white, middle color=orange!60!brown!60]
		([xshift=-5pt]frame.south west) -- ([xshift=5pt]frame.south east);
	},
	#1
}

\newcommand{\highlightbox}[1]{%
	\begin{fancyhighlightorange}#1\end{fancyhighlightorange}%
}

% -----------------------------------------------------------
% DIAGRAM SCALED PLACEMENT
% -----------------------------------------------------------
\providecommand{\myscalebelow}{0.7}
\providecommand{\myscaleright}{0.55}

% Required in the preamble
\usepackage{xcolor}
\usepackage{ifthen}

\newif\ifPlacePNGmissing

\newcommand{\PlacePNG}[4]{%
	\PlacePNGmissingfalse
	%
	% Check whether any argument is empty
	\if\relax\detokenize{#1}\relax\PlacePNGmissingtrue\fi
	\if\relax\detokenize{#2}\relax\PlacePNGmissingtrue\fi
	\if\relax\detokenize{#3}\relax\PlacePNGmissingtrue\fi
	\if\relax\detokenize{#4}\relax\PlacePNGmissingtrue\fi
	%
	\ifPlacePNGmissing
	\par\noindent
	{\color{orange}\bfseries A FIGURE COMES HERE}
	\par\vspace{1em}
	\else
	\ifthenelse{\equal{#3}{below}}%
	{\def\myscale{\myscalebelow}}%
	{\def\myscale{\myscaleright}}%
	%
	\ifthenelse{\equal{#3}{below}}{%
		\par\noindent
		\begin{minipage}{\linewidth}
			\centering
			\includegraphics[width=\myscale\linewidth]{#2}
		\end{minipage}
		\par\vspace{1em}
	}{%
		\ifthenelse{\equal{#3}{right}}{%
			\noindent
			\begin{minipage}[t]{0.57\linewidth}
				\vspace{0pt}#4
			\end{minipage}
			\hfill
			\begin{minipage}[t]{0.4\linewidth}
				\vspace{0pt}
				\includegraphics[width=\myscale\linewidth]{#2}
			\end{minipage}
		}{%
			\textbf{Error: position must be "below" or "right".}
		}%
	}%
	\fi
}




% Define the HypThmBox command.
% This command takes two arguments: the Hypothesis text and the Thesis text.
\newcommand{\HypThmBox}[2]{%
	\begin{tcolorbox}[
		enhanced,
		colback=gray!4!white,
		colframe=black,
		boxrule=0.8pt,
		arc=4pt,
		left=2mm,right=2mm,top=2mm,bottom=2mm,
		width=\textwidth,
		sidebyside,
		sidebyside gap=4mm,
		lefthand width=.48\textwidth,
		righthand width=.48\textwidth,
		separator sign={\MySeparator},
		before upper={\textit{\textbf{Hypothesis}}\\[8pt]},
		before lower={\textit{\textbf{Thesis}}\\[8pt]}
		]
		#1
		\tcblower
		#2
	\end{tcolorbox}%
}


% -----------------------------------------------------------
% EXERCISE AND TASKS
% -----------------------------------------------------------
\usepackage{exercise}
\usepackage{tasks}
\setlength{\columnsep}{0.7cm}

% -----------------------------------------------------------
% CUSTOM COMMANDS
% -----------------------------------------------------------
\newcommand{\overbar}[1]{\mkern 1.5mu\overline{\mkern-1.5mu#1\mkern-1.5mu}\mkern 1.5mu}
\newcommand{\norm}[1]{\left\lVert#1\right\rVert}
\newcommand{\vect}[1]{\underline{#1}}
\newcommand{\binary}[2]{$#1$-$#2$}
\newcommand{\dom}{{\rm dom\ }}
\newcommand{\ran}{{\rm ran\ }}
\newcommand{\LCM}{{\rm LCM\ }}
\newcommand{\HCF}{{\rm HCF\ }}
\newcommand{\ep}{\epsilon}
\newcommand{\un}{\mathscr{U}}
\newcommand{\sm}{\backslash}
\newcommand{\lb}{\vspace{0.1cm}}
\newcommand{\Lb}{\vspace{0.2cm}}
\newcommand{\Bb}{\vspace{0.4cm}}
\newcommand{\HRule}[1]{\rule{\linewidth}{#1}}
\newcommand{\und}[1]{\underline{\smash{#1}}}
\newcommand{\der}{\text{d}}
\newcommand{\naturals}{\mathbb{N}}
\newcommand{\integer}{\mathbb{Z}}
\newcommand{\rational}{\mathbb{Q}}
\newcommand{\real}{\mathbb{R}}
\newcommand{\complex}{\mathbb{C}}
\newcommand{\blankpage}{\newpage\thispagestyle{empty}\mbox{}\newpage}
\newcommand{\divides}{\bigm|}
\newcommand{\ndivides}{%
	\mathrel{\mkern.5mu\ooalign{\hidewidth$\big|$\hidewidth\cr$\nmid$\cr}}%
}

% -----------------------------------------------------------
% TIKZ PICTURE SCALING
% -----------------------------------------------------------
\makeatletter
\usepackage{environ}
\newsavebox{\measure@tikzpicture}
\NewEnviron{scaletikzpicturetowidth}[1]{%
	\def\tikz@width{#1}%
	\def\tikzscale{1}\begin{lrbox}{\measure@tikzpicture}%
		\BODY
	\end{lrbox}%
	\pgfmathparse{#1/\wd\measure@tikzpicture}%
	\edef\tikzscale{\pgfmathresult}%
	\BODY
}
\makeatother





